\documentclass[11pt]{article}
\usepackage{amsmath, amssymb, pgfplots}
\pgfplotsset{compat=1.18}
\newcommand{\R}{\mathbb{R}}
\newcommand{\dd}[2]{\frac{d #1}{d #2}}
\DeclareMathOperator{\sgn}{sgn}
\title{course Homework 5: Continuity}
\author{Prof.~Example}

\begin{document}
\maketitle

\section{Limits and continuity} % comments are not text
Let $f$ be defined on all of $\R$. A function is \textbf{continuous} at $a$ if
\[
  \lim_{x \to a} f(x) = f(a).
\]
We write $\dd{y}{x}$ for the derivative, and 50\% of the points come from part (b).

\begin{enumerate}
  \item Find $\lim_{x \to 2} \frac{x^2 - 4}{x - 2}$.
  \item Is $g$ continuous at $x = 0$?
  \begin{equation*}
    g(x) = \begin{cases} x^2 & x < 0 \\ 2x + 1 & x \ge 0 \end{cases}
  \end{equation*}
  \begin{enumerate}
    \item Find the left limit.
    \item Find the right limit.
  \end{enumerate}
  \item Simplify:
\end{enumerate}

\begin{align*}
  (x+1)^2 &= x^2 + 2x + 1 \\
          &= (x+1)(x+1) \label{eq:square}
\end{align*}

\subsection*{A table}
\begin{tabular}{|c|c|}
\hline
$x$ & $f(x)$ \\ \hline
1 & 3 \\
2 & 4 \\ \hline
\end{tabular}

\begin{tikzpicture}
\begin{axis}[xmin=-4, xmax=4, ymin=-2, ymax=2, axis lines=middle]
  \addplot[domain=-4:4, samples=100] {sin(deg(x))};
  \addplot[only marks, mark=*] coordinates {(0,0) (1.5708,1)};
\end{axis}
\end{tikzpicture}

\begin{theorem}[Intermediate Value Theorem]
If $f$ is continuous on $[a, b]$ and $N$ is between $f(a)$ and $f(b)$, then $f(c) = N$ for some $c$.
\end{theorem}

\newpage
\section{Graphs}
\includegraphics[width=3in]{parabola.png}
``Quoted'' text --- with a dash, and a \href{https://example.com}{link}.
\end{document}
